\section{Related Work}

\subsection{Indirect prompt injection and RAG poisoning}
Indirect prompt injection exploits the fact that language models process
external content as well as user instructions. Greshake et al. demonstrate
that retrieved or otherwise externally supplied text can steer
LLM-integrated applications, including by changing their behavior and
triggering unintended actions \cite{greshake2023indirect}. Yi et al.
introduce BIPIA as a benchmark for measuring this risk across tasks in which
malicious instructions are embedded in ordinary external content
\cite{yi2023benchmarking}. Perez and Ribeiro show that instructions such as
``ignore previous instructions'' can hijack a model's goal, which is the kind
of phrasing that instruction-pattern matching targets~\cite{perez2022ignore}.
Liu et al.\ formalize prompt injection and benchmark attacks and defenses
\cite{liu2024formalizing}. These works motivate
treating retrieved documents as untrusted inputs, while also showing that
document-level screening is only one possible defense layer.

RAG knowledge poisoning targets the contents of the retrieval database.
PoisonedRAG formulates attacks that insert a small number of malicious texts
to steer answers to attacker-chosen targets \cite{zou2024poisonedrag}. Zhong
et al. show that a few adversarial passages inserted into a corpus can be
retrieved for many queries~\cite{zhong2023poisoning}; AgentPoison poisons the
memory or knowledge base of LLM agents~\cite{chen2024agentpoison}; and blocker
documents make a RAG system withhold an answer~\cite{shafran2025jamming}. We did
not evaluate against these retrieval-side attacks. Chen
et al. study detection and removal of indirect injections in retrieved
documents and provide a benchmark for both tasks \cite{chen2025detectremove}.
Their work is close to PDS in its document-level focus, but evaluates
instruction detection and removal models; PDS instead compares existing
pattern, embedding, perplexity, and learned-fusion signals under a fixed
false-alarm protocol.

\subsection{Prompt-injection datasets and detectors}
The deepset prompt-injections dataset provides labeled benign and injection
examples and is one source of training data for open prompt-injection
classifiers \cite{deepset2024dataset}. Protect AI's DeBERTa-v3-base
prompt-injection model is a published classifier trained on a mixture of
public datasets and crafted injections; its model card describes English
input and notes limitations on jailbreak and non-English detection
\cite{protectai2024deberta}. These resources provide useful comparison
points, but their reported performance depends on their training data and
evaluation distributions. PDS evaluates the published classifier alongside
its own signals, one of which is a sentence classifier fitted on the training
role of our data.

Other defense layers act on the model rather than on the document.
Spotlighting marks untrusted text so that the model can tell it apart from
instructions~\cite{hines2024spotlighting}, and StruQ separates instructions from
data in a structured query and trains the model to follow only the instruction
channel~\cite{chen2025struq}. PDS screens documents before they reach the
model, so it can be combined with these.

Perplexity filtering is another lightweight signal. Alon and Kamfonas
evaluate perplexity for detecting language-model-generated adversarial
attacks and report that false positives are a significant challenge for
plain perplexity filtering \cite{alon2023perplexity}. Jain et al. likewise
evaluate a perplexity filter as a baseline defense against adversarial attacks
on aligned language models~\cite{jain2023baseline}. This motivates
measuring fluency as one component of a detector rather than treating it as
a sufficient standalone decision rule.

\subsection{Embedding comparison and calibration}
CED compares text embeddings to identify out-of-distribution or hallucinated
text, while emphasizing that these categories can overlap in embedding space
\cite{lee2024ced}. This is relevant to PDS's embedding-anomaly signal and
its observed sensitivity to domain shift: distance from a clean reference
corpus is not itself evidence of malicious intent.

Conformal prediction provides distribution-free, finite-sample uncertainty
sets under exchangeability assumptions and can be applied on top of a
pretrained model \cite{angelopoulos2022conformal}. Such guarantees concern
the specified conformal procedure and assumptions; they do not follow from
selecting a threshold at an empirical false-alarm budget alone. PDS therefore
reports threshold-based performance under its fixed protocol and treats
calibration as a separate design question.

\subsection{Evaluation practice in security}
Arp et al. review machine-learning papers in security and list recurring
pitfalls, among them sampling bias, data snooping, spurious correlations and the
base-rate fallacy~\cite{arp2022dos}. Axelsson showed for intrusion detection that
the base rate of attacks, not the detector's accuracy alone, decides how many
alerts are false~\cite{axelsson2000base}. A related warning comes from natural
language inference, where hypothesis-only baselines, which never see the
premise, performed well above the majority class on a number of datasets and so
exposed artefacts in how the data were built~\cite{poliak2018hypothesis}.
Section~\ref{sec:protocol} adopts grouped splits, a held-out role scored once and
precision at low prevalence for these reasons, and
Section~\ref{sec:learned} applies a baseline of the same kind.

\subsection{Positioning of this work}
PDS does not claim a new detection principle: the signals are existing ones,
and the learned parts are a linear sentence classifier and a logistic-regression
fusion. Its contribution is a failure-and-latency comparison of existing signals
and their combinations under a fixed false-alarm budget, with attention to source shift, adversarial examples,
and operational trade-offs. The related work establishes the threat and
candidate signals; our evaluation asks which signals complement one another,
what their errors look like, and what their computational cost buys.